\documentclass[10pt,a4paper]{amsart}
\usepackage[margin=19mm]{geometry}
\usepackage{amsmath,amssymb,mathtools}
\usepackage[scr=rsfs,cal=euler]{mathalfa}
\usepackage{xfrac}
\usepackage[unicode,hidelinks,bookmarks=false]{hyperref}
\newcommand{\ie}{i.e.\ }
\newcommand{\cf}{cf.\ }
\newcommand{\resp}{respectively\ }
\newcommand{\Subsection}{\subsection}
\providecommand{\nobreakdash}{\nobreak\hbox{-}}
\setlength{\emergencystretch}{2em}
\multlinegap0pt
\usepackage{amssymb}
\usepackage{tikz}
\newtheorem{prop}{Proposition}
\newtheorem{lemma}{Lemma}
\usepackage{cleveref}
\crefname{prop}{Proposition}{Propositions}
\crefname{lemma}{Lemma}{Lemmas}
\DeclareMathOperator{\rk}{rk}
\title{\bf{Piecewise isometry groups of Euclidean tessellations}}
\author{Robert Bieri, Alex Feingold, Daniel Studenmund}
\date{Source: arXiv:2607.28893v1 (CC BY 4.0)}
\begin{document}
\maketitle

\noindent\emph{Source and licence.} \bf{Piecewise isometry groups of Euclidean tessellations}. Version arXiv:2607.28893v1 (CC BY 4.0), distributed by arXiv under the Creative Commons Attribution 4.0 International licence, http://creativecommons.org/licenses/by/4.0/, as recorded in the arXiv licence metadata for this exact version and shown on the abstract page https://arxiv.org/abs/2607.28893v1. Full licence notice: \texttt{licenses/arxiv-2607.28893-CC-BY-4.0.txt}.\par\smallskip

\noindent\textit{Editorial scope.} Counted unit: the article's lemma lem:irred-form (source0.tex lines 932--941). The article's complete original proof of it is delivered with it (source0.tex lines 943--985, one \texttt{proof} environment of 43 lines). No mathematics is written, repaired or re-derived: the statement and the proof are the article's own lines, in the article's own notation, and no retained line is substituted or re-flowed.  Retained uncounted as the source-local context the proof needs: lines 853--864.  Nothing was omitted from any retained span, so the package carries no omission record.  No retained line contains a citation, so the wrapper carries no bibliography and no bibliography entry.  No retained line contains a figure environment, an image-inclusion command or drawing code, so no image is delivered and the package declares no asset dependency.  Editorial formatting only, in addition: an AMS \texttt{amsart} preamble that restates the article's own declarations and macros used by the retained lines, namely the environments \texttt{prop}, \texttt{lemma} on separate counters, together with the standard packages those lines need.  The retained environments print in this document as Proposition 1, Lemma 1 in order of appearance, whereas the article prints the same items with the numbering of its own full document.  The complete native source of the article as distributed in the arXiv e-print for this exact version is bundled unchanged as \texttt{sources/206446/source0.tex}.
\begin{prop} \label{prop:piececut} Suppose $P$ is a nonempty
  convex polyhedral set.  Fix $\xi\in \Xi$ and a half-space
  $M\in \xi$ such that $P\doublecap M$ is nonempty.
  \begin{enumerate}
  \item \label{unboundcutpart} If $P$ is $\xi$-bounded, then  
    $\rk(P\doublecap M) = \rk(P)$.
  \item \label{boundcutpart} If $P$ is $\xi$-unbounded and 
    $\iota(\xi)$-bounded, then  $\rk(P\doublecap M) < \rk(P)$.
  \item \label{doubleunboundcutpart} If $P$ is unbounded with respect
    to both $\xi$ and $\iota(\xi)$, then $\rk(P\doublecap M) = \rk(P)$.
  \end{enumerate}
\end{prop}

\begin{lemma} \label{lem:irred-form} For a nonempty convex polyhedral
  set $P$, the following are equivalent:
  \begin{enumerate}
  \item \label{cond-irred} $P$ is irreducible.
  \item \label{cond-bound} $P$ is thin and for each $\xi \in \Xi$, $P$
    is $\xi$-bounded or $\iota(\xi)$-bounded.
  \item $P$ is thin and $L(P)$ is a closed cell in the spherical
    complex $L(\Delta)$.
  \end{enumerate}
\end{lemma}

\begin{proof}
  To see that irreducibility implies the second condition, suppose $P$
  is irreducible. First suppose there exists $\xi \in \Xi$ such that
  $P$ is both $\xi$-unbounded and $\iota(\xi)$-unbounded. Then for any
  $M\in \xi$, both $P \cap M$ and $P\cap \iota(M)$ have rank equal to
  the rank of $P$ by \cref{prop:piececut}, while
  $P = (P\cap M) \cup (P\cap \iota(M))$ is a finite decomposition,
  contradicting irreducibility of $P$.
  
  Now suppose $P\subset M\cap M'$ where $M\in \xi$ and
  $M'\in \iota(\xi)$, and there is $M''\in\xi$ such that both
  $P\doublecap M''$ and $P\doublecap \iota(M'')$ are nonempty. Then
  both $P \cap M''$ and $P\cap \iota(M'')$ have rank equal to the rank
  of $P$ by \cref{prop:piececut}, while
  $P = (P\cap M'') \cup (P\cap \iota(M''))$ is a finite decomposition,
  contradicting irreducibility of $P$.
  
  To see that the second condition irreducibility, suppose $P$ is
  reducible, so that there is a finite decomposition
  $P = (P\cap M) \cup (P\cap \iota(M))$ for some $M\in \xi$, where
  both $P\cap M$ and $P\cap \iota(M)$ have the same rank as
  $P$. Suppose that $P$ is $\xi$-bounded or
  $\iota(\xi)$-bounded. By \cref{prop:piececut}, $P$ must be both
  $\xi$-bounded and $\iota(\xi)$-bounded. But since both
  $P\cap M$ and $P\cap \iota(M)$ are nonempty, $P$ cannot be thin.

  To finish the proof, we assume $P$ is thin and show that $P$ is
  $\xi$-bounded or $\iota(\xi)$-bounded for each $\xi\in \Xi$ if and
  only if its limit set is a closed cell in $L(\Delta)$.  Let $\Xi_P$
  be the set of $\xi\in \Xi$ for which $P$ is $\xi$-bounded. Consider
  $\Xi_P \cap \iota(\Xi_P)$, the set of $\xi$ for which $P$ is both
  $\xi$-bounded and $\iota(\xi)$-bounded. Consider the subsphere
  \[
    S = \bigcap_{\xi \in \Xi_P\cap \iota(\Xi_P)} L(\xi).
  \]
  Now, $L(P)$ is the intersection of $S$ with a closed half-space for
  each $\beta \in \Xi$ for which $P$ is $\beta$-bounded but not
  $\iota(\beta)$-bounded. The desired result now follows because the
  resulting set is a closed cell in $S$ precisely when this
  intersection includes exactly one half-space
  $M\in \beta\cup \iota(\beta)$ for each
  $\beta \in \Xi - (\Xi_P \cap \iota(\Xi_P))$.
\end{proof}

\end{document}
